\documentclass[11pt]{article}
\usepackage[margin=1in]{geometry}
\usepackage{graphicx}
\usepackage{booktabs}
\usepackage{amsmath,amssymb}
\usepackage{microtype}
\usepackage{hyperref}
\usepackage{xcolor}
\usepackage{enumitem}
\usepackage[T1]{fontenc}
\usepackage[utf8]{inputenc}
\usepackage{float}

\graphicspath{{../results/figures/}}

\definecolor{note}{RGB}{60,80,140}
\definecolor{newc}{RGB}{20,120,40}
\newcommand{\note}[1]{{\color{note}\textit{#1}}}
\newcommand{\honest}[1]{\textcolor{red!70!black}{\textbf{Honesty note:} #1}}
\newcommand{\newin}[1]{\textcolor{newc}{\textbf{New in v2:} #1}}

\title{An Explainer (v2):\\
Ablation-First Trajectory Analysis of Alveolar Epithelial\\
Robustness Collapse in Lethal COVID-19\\
\large{with Harmony-Corrected Batch Analysis and\\
Independent Cross-Cohort Replication}}
\author{A guide for computational biologists, computer scientists, and biologists}
\date{\today}

\begin{document}
\maketitle

\begin{abstract}
\noindent This explainer walks through the revised (v2) analysis of the Columbia/NYP
COVID-19 Lung Atlas (SCP1219) end-to-end, with three new pieces of work layered onto
the original primary-cohort analysis:
\begin{enumerate}
\item a \textbf{fix for the Harmony batch-correction step} that had silently degenerated
      Ablation 3 in the original paper;
\item an \textbf{independent cross-cohort replication} on 89{,}736 AT1/AT2 cells from
      618 donors across 35 datasets (cellxgene census, HLCA-extended + Krasnow, with
      SCP1219 excluded); and
\item a \textbf{gene-program audit} that resolved a previously open HVG-dropout concern.
\end{enumerate}
The biology, statistical framework, ablation logic, and reading-the-plots guidance from
the original explainer are retained --- and are annotated here to show which claims
strengthen, which change direction, and which are retracted under the new evidence.
The v2 bottom line is unchanged in essence but far better supported: robustness collapse
(dispersion + donor-level pseudotime displacement) transfers across cohorts; directional
centroid displacement and marker-threshold DATP enrichment do not.
\end{abstract}

\tableofcontents
\bigskip

\section{What changed between v1 and v2}

The original explainer and paper reported three pre-registered tests of
``epithelial robustness collapse'' on a single cohort (SCP1219, Melms 2021). Two passed
(dispersion, pseudotime shift) and one failed (centroid displacement). The paper's
credibility rested on the fact that those tests were genuinely pre-registered and
honestly reported.

In v2 we do three additional pieces of work that either strengthen or stress-test those
primary-cohort conclusions:

\begin{description}[style=nextline]
\item[(i) Harmony fix.] The primary paper listed as a limitation that
      ``Harmony failed at runtime; Ablation~3 cannot distinguish corrected from
      uncorrected.'' The root cause was a shape-assignment bug in the wrapper that
      assumed harmonypy~$<$0.2 semantics; the current PyTorch backend returns
      \texttt{Z\_corr} already in (cells,~PCs) orientation, and the old transpose
      produced a degenerate matrix. We added an orientation-detecting check. With the
      fix, Harmony converges cleanly and the Ablation~3 comparison is a real test.
      The corrected branch \emph{amplifies} the pseudotime shift from $+0.0483$ to
      $+0.1246$ (2.6$\times$) while leaving the dispersion and displacement statistics
      unchanged.
\item[(ii) Gene-program audit.] The primary paper flagged GPX1 as absent from the HVG
      set and noted that other programs might have analogous silent dropouts. We
      confirmed that GPX1 is absent from the SCP1219 reference gene-name file in the
      first place (not a downstream filtering artefact); the other 83 program genes are
      all present upstream. This retracts the ``other programs may have analogous
      dropouts'' hedge.
\item[(iii) Independent cross-cohort replication.] We queried cellxgene census
      (2025-11-08) for all human lung pulmonary alveolar type~1/2 cells under
      \texttt{disease==COVID-19} or \texttt{disease==normal}, strictly excluding SCP1219
      by \texttt{dataset\_id}. After a per-donor 200-cell cap (RNG seed 42) the
      replication set contains 89{,}736 cells (4{,}441 COVID / 85{,}295 normal) from
      618 donors and 35 datasets. We score the same 8 gene programs on this cohort and
      test the same three robustness-collapse claims.
\end{description}

\newin{this explainer has a new \S\ref{sec:harmony-fix} on the Harmony bug and fix, a
new \S\ref{sec:replication} on the replication, and annotations throughout the biology
and methods sections noting where conclusions changed.}

\section{The biology, in plain terms}

\subsection{What is the alveolar epithelium?}

Deep in the lung, roughly 480 million tiny air sacs (\emph{alveoli}) are where you
actually exchange O\textsubscript{2} and CO\textsubscript{2} with your blood. The wall
of each alveolus is lined by a two-cell epithelium:

\begin{description}[style=nextline]
\item[AT1 (alveolar type I) cells] Extraordinarily flat; they cover about 95\% of the
      alveolar surface. They are the \emph{gas-exchange barrier}.
\item[AT2 (alveolar type II) cells] Cuboidal; they secrete \emph{surfactant} and act as
      the local \emph{stem cells}: when AT1 cells are destroyed, AT2 cells divide and
      differentiate into new AT1 cells through a transitional state marked by the genes
      \textbf{KRT8}, \textbf{CLDN4}, and \textbf{SFN}. These transitional cells are
      often called DATPs (damage-associated transient progenitors).
\end{description}

\subsection{What goes wrong in lethal COVID-19?}

SARS-CoV-2 enters AT2 cells via the ACE2 receptor. The virus is infecting the very
cells that the lung depends on both for surfactant and for regeneration. On top of that,
the immune response itself damages the epithelium (inflammatory ``collateral damage'').
Autopsy studies describe diffuse alveolar damage: AT1 cells missing, edema, and
abnormal cells that are neither cleanly AT1 nor cleanly AT2 --- apparently caught
mid-repair. The core question: does the alveolar epithelium fail as a \emph{coherent
block} (all cells move to a new damaged state), or does it fail as a \emph{graded,
continuous process} (cells fan out along a spectrum of states)?

\subsection{What is already known --- and therefore what this project is \emph{not}}

Melms 2021 (the SCP1219 paper itself), Delorey 2021, and Bharat 2020 have already
described COVID-19 alveolar injury in detail, including DATP accumulation, AT2 stress
signatures, and donor-to-donor heterogeneity. This project is \textbf{not} a new
biology discovery; it is a \textbf{methods} contribution: what happens when you stress
a trajectory-analysis reading of this dataset along ten pre-registered axes, and do the
robustness-collapse signatures survive an independent 618-donor replication?

\section{The data}

\subsection{Primary cohort (SCP1219)}

Post-mortem lung tissue from 20 COVID-19 and 7 healthy organ donors;
116{,}313 snRNA-seq nuclei. After QC and restriction to alveolar epithelium we keep
22{,}128 cells (9{,}608 AT1, 11{,}341 AT2, 1{,}179 ECM-high/transitional). Per-donor
counts range 62--2{,}850; 10 of 27 donors have $<$300 cells.

\note{snRNA-seq rather than scRNA-seq because whole cells don't survive post-mortem
freeze/thaw but nuclei do. Trade-off: captures more unspliced/nascent transcript.}

\subsection{Replication cohort (cellxgene census)}\label{sec:replication-data}

We deliberately held Delorey (GSE171524) and Wendisch (PRJNA705687) in mind as the
ideal full-transcriptome replication targets, but neither is hosted in cellxgene census
as of 2025-11-08. As a practical alternative we took all human lung AT1/AT2 cells in
the census under \texttt{disease==COVID-19} or \texttt{disease==normal} with
\texttt{dataset\_id !=} SCP1219. This is drawn overwhelmingly from the HLCA-extended
integration and the Krasnow COVID-infection study, and it spans 35 datasets and
618 donors.

We deliberately restricted the gene axis to the 83 genes referenced by our 8 gene
programs plus 20 canonical markers and SARS-CoV-2 entry genes (84 unique symbols total),
keeping the download to a tractable size. The consequence: the replication cohort is
rich enough to rescore the gene programs and run a marker-based DATP assay, but
\emph{not} rich enough to build a new HVG/PCA/diffusion manifold and repeat diffusion
pseudotime from scratch. This is disclosed as a remaining limitation.

\subsubsection*{Class-balance asymmetry and the per-donor cap}

Before capping, the replication cohort is $\approx$15:1 normal:COVID. This reflects
that cellxgene census is mostly healthy atlas data with a few COVID studies mixed in.
A cell-level test would be dominated by whichever few large normal donors happen to
have extreme scores; we mitigate by capping each donor to $\leq$200 cells with seed~42
(89{,}736 cells post-cap) and by always reporting a donor-level pseudobulk test
alongside the cell-level one.

\section{Pipeline --- what each step computes and why}

Versions: scanpy 1.12.1, anndata 0.12.10, harmonypy 0.2.0, palantir 1.4.4,
Python 3.12.3, RNG seed 42.

\begin{enumerate}
\item \textbf{QC.} $\geq$500 UMI, $\geq$250 genes, $\leq$20\% mitochondrial, doublet
      filter. QC tables in \texttt{results/tables/tableS2\_qc\_*.csv}.
\item \textbf{Normalization.} \texttt{sc.pp.normalize\_total(target=1e4)} then
      \texttt{log1p}. \texttt{.raw} captures the full-gene log-normalized matrix for
      gene-program scoring.
\item \textbf{HVG selection.} 3{,}000 Seurat-v3 HVGs.
\item \textbf{Scale \& PCA.} scale to unit variance (clip at 10), 50-component PCA.
\item \textbf{Batch correction.} Harmony on \texttt{donor\_id}
      (see \S\ref{sec:harmony-fix} for the v2 fix).
\item \textbf{Neighbor graph.} $k$-NN with $k{=}15$ on PCA-50 (uncorrected by default;
      Harmony branch is an ablation).
\item \textbf{Embeddings.} UMAP (visualization) and 15-component diffusion map
      (trajectory).
\item \textbf{Pseudotime.} Diffusion pseudotime rooted at the healthy-AT2 centroid,
      normalized to $[0,1]$.
\item \textbf{Gene-program scoring.} \texttt{sc.tl.score\_genes} against \texttt{.raw}
      for 8 programs (AT1 identity, AT2 identity, interferon response, NF-$\kappa$B,
      oxidative stress, apoptosis, senescence, AT2$\to$AT1 differentiation).
\item \textbf{Robustness tests.}
      (a)~one-sided Mann--Whitney $U$ on PCA-50 distances (centroid displacement);
      (b)~Levene's test (dispersion);
      (c)~1000-permutation test on median pseudotime difference;
      (d)~DATP (KRT8$^+$/CLDN4$^+$) chi-square.
\end{enumerate}

\section{The Harmony fix}\label{sec:harmony-fix}

\subsection{What went wrong}

The v1 paper included the Harmony branch of Ablation~3 but the output was identical to
the uncorrected branch because of a wrapper bug. Specifically, the wrapper assigned

\begin{verbatim}
adata.obsm["X_pca_harmony"] = ho.Z_corr.T
\end{verbatim}

This followed the harmonypy~$<$~0.2 convention (Z\_corr in PCs$\times$cells). The 0.2
release switched to a PyTorch backend that returns Z\_corr in cells$\times$PCs.
The transpose then produced a (50, 22128) matrix which either failed assignment or was
silently matched against the wrong axis downstream, giving a degenerate embedding that
mimicked the uncorrected PCA.

\subsection{The fix}

We replaced the unconditional transpose with an orientation check
(\texttt{src/embedding.py} \texttt{run\_harmony}):

\begin{verbatim}
Z = ho.Z_corr
if Z.shape[0] != adata.n_obs and Z.shape[1] == adata.n_obs:
    Z = Z.T
adata.obsm["X_pca_harmony"] = Z
\end{verbatim}

Harmony now converges in 3 iterations.

\subsection{What the working Ablation 3 shows}

\begin{table}[H]
\centering
\begin{tabular}{lrrr}
\toprule
Branch & MPD (COVID$-$control) & $\rho$ & $r$ \\
\midrule
Uncorrected PCA & $+0.0483$ & 0.548 & 0.141 \\
Harmony-corrected PCA & $+0.1246$ & 0.452 & 0.141 \\
\bottomrule
\end{tabular}
\caption{Ablation 3, v2. Harmony $2.6\times$ amplifies the pseudotime shift while
leaving the dispersion/displacement statistics unchanged.}
\end{table}

This is a qualitative change in how we should interpret the v1 paper. Previously the
v1 paper could not say what batch correction did to the signal --- it was a genuine
open limitation. The v2 evidence is that donor-level batch correction strengthens the
primary conclusion. That is the direction of evidence you hope to see, but cannot
assume: Harmony could equally have flattened the shift.

\section{Reading the primary-cohort results (unchanged)}

The three pre-registered tests and their outcomes on SCP1219 are unchanged between
v1 and v2 --- they depend only on the primary pipeline, not on the Harmony fix or the
replication.

\subsection{Centroid displacement (fails)}

One-sided Mann--Whitney $U$ in PCA-50 (COVID $>$ Healthy): $p=1.0$; COVID median
lower than healthy; $r=0.14$ opposite to the alternative. The pre-specified core test
did not distinguish conditions directionally.

\subsection{Dispersion (supported)}

Within-group distance to own centroid, PCA-50. COVID variance $33.28$; healthy $14.58$
(2.3$\times$). Levene $=276.5$, $p<10^{-30}$. COVID cells are \emph{spread out}, not
\emph{pushed away}.

\subsection{Pseudotime enrichment (supported)}

Control median $0.063$; COVID median $0.111$; observed median difference $+0.0483$.
1000-permutation null: mean $6.0\times10^{-5}$, std $1.24\times10^{-3}$. $p=0.000999$.
Leave-one-donor-out preserves direction in all 27 iterations.

\section{The ablation battery, annotated with v2}

All ten ablations from the primary paper are unchanged except for Ablation~3.
Table~\ref{tab:ablations-v2} presents the complete set using the v2 Harmony numbers.

\begin{table}[H]
\scriptsize
\centering
\begin{tabular}{lrrrl}
\toprule
Ablation & MPD & $\rho$ & $r$ & Notes \\
\midrule
1 AT2-only & $+0.133$ & 0.33 & $-0.29$ & sign flip on $r$ \\
2 UMAP / diffmap & $+0.048$ & 0.55 & $+0.14$ & unchanged \\
3 Harmony corrected & $+0.125$ & 0.45 & $+0.14$ & \textbf{v2 --- now functional} \\
3 Uncorrected & $+0.048$ & 0.55 & $+0.14$ & baseline \\
4 No apoptosis genes & $+0.074$ & 0.50 & $+0.14$ & slight up \\
5 4 roots & $-0.037$ to $+0.048$ & $-0.10$ to 0.64 & $+0.14$ & root-dependent \\
6 Broader epithelial & $-0.008$ & 0.83 & $+0.25$ & shift abolished \\
7 LOO (27) & 0.028--0.202 & $-0.17$--0.74 & per-iter & direction 27/27 \\
8 Curated-only & $+0.048$ & 1.00 & $+0.14$ & stable \\
8 MSigDB-only & $+0.048$ & $-0.80$ & $+0.14$ & ordering inverted \\
9 Palantir & $+0.028$ & --- & --- & alt pseudotime algorithm \\
10 Permutation & $p=0.001$ & --- & --- & 1000 shuffles \\
\bottomrule
\end{tabular}
\caption{Ten-way ablation summary, v2.}\label{tab:ablations-v2}
\end{table}

\section{Independent replication}\label{sec:replication}

\subsection{Injury composite surrogate}

Because the replication cohort has only 84 genes, we cannot rebuild the manifold.
As a surrogate we define an \emph{injury composite} $=$ the mean of six
non-homeostatic program scores (interferon, NF-$\kappa$B, oxidative stress, apoptosis,
senescence, AT2$\to$AT1). Higher values $\Rightarrow$ more injury-program activation.

\subsection{Dispersion --- replicates strongly}

\begin{table}[H]
\centering
\begin{tabular}{lrr}
\toprule
Metric & COVID & Normal \\
\midrule
Injury composite variance & 0.0505 & 0.0306 \\
Variance ratio COVID/Normal & \multicolumn{2}{r}{\textbf{1.65$\times$}} \\
Levene statistic & \multicolumn{2}{r}{626.18} \\
Levene $p$ & \multicolumn{2}{r}{$1.0\times 10^{-137}$} \\
\bottomrule
\end{tabular}
\caption{Replication dispersion test. The variance ratio drops from 2.3$\times$
in the primary cohort to 1.65$\times$ in the 618-donor replication, consistent with
cross-study heterogeneity dampening effect sizes without eliminating them.}
\end{table}

\subsection{Pseudotime surrogate --- cell-level fails, donor-level replicates}

One-sided Mann--Whitney (COVID $>$ normal) on cell-level injury composite: $p=1.0$.
Donor-level pseudobulk (43 COVID donors vs 575 normal donors): $p=1.98\times10^{-3}$
\emph{in the predicted direction}.

This inversion is informative. In a class-imbalanced heterogeneous atlas,
cell-level tests are driven by the largest donors of the majority class. Donor-level
pseudobulk is the right unit of analysis for cross-study tests, and it replicates.

\subsection{KRT8$^+$/CLDN4$^+$ DATP enrichment --- does NOT replicate}

\honest{in the primary cohort we saw $3.1\times$ enrichment of KRT8$^+$/CLDN4$^+$ cells
in COVID. In the replication cohort we see the \emph{opposite} direction: COVID
$22.9\%$ vs normal $47.6\%$, fold $=0.48\times$ ($\chi^2\ p<10^{-200}$). The most
plausible explanation is that a simple expression-threshold cutoff is not portable
across 35 datasets with different dissociation protocols, sequencing depths, and
normalization choices. The DATP \emph{biology} is well-supported by manifold-based
identifications in Kobayashi/Strunz/Adams/Habermann 2020 and by the Melms paper itself,
but the marker-threshold chi-square assay we used should be \textbf{retracted as a
portable population-level metric}.}

\section{Composite assessment (v2)}

\begin{table}[H]
\centering
\small
\begin{tabular}{lll}
\toprule
Criterion & Primary (SCP1219) & Replication (census) \\
\midrule
Centroid displacement & fails ($p=1.0$) & fails cell-level ($p=1.0$) \\
Dispersion (Levene) & supported ($p\approx 0$) & supported ($p\approx10^{-137}$) \\
Pseudotime enrichment & supported ($p=0.001$) & supported donor-level ($p=1.98\text{e-}3$) \\
Harmony-corrected shift & --- & MPD $\times 2.6$ \\
KRT8$^+$/CLDN4$^+$ DATPs & $3.1\times$ enriched & $0.48\times$ (retracted) \\
\bottomrule
\end{tabular}
\caption{v2 assessment across both cohorts.}
\end{table}

\section{What the v2 evidence supports}

\begin{enumerate}
\item Robustness collapse (dispersion + donor-level pseudotime displacement)
      \emph{transfers} across cohorts.
\item Harmony batch correction \emph{amplifies} the pseudotime shift by $\approx$2.6$\times$.
\item Dispersion (Levene) is a portable, cross-study-robust signal of
      robustness collapse.
\item Donor-level pseudobulk Mann--Whitney is the correct unit of analysis for
      cross-study replication in imbalanced atlases.
\end{enumerate}

\section{What the v2 evidence does NOT support}

\begin{enumerate}
\item Directional mono-vector centroid displacement --- fails in both cohorts.
\item A specific IFN$\to$NF-$\kappa$B$\to$repair$\to$apoptosis cascade
      (ordering inverts under MSigDB or COVID-extreme root).
\item Marker-threshold KRT8$^+$/CLDN4$^+$ DATP population enrichment as a portable
      cross-cohort statistic.
\item Any claim about elapsed time --- pseudotime is a manifold ordering.
\end{enumerate}

\section{Remaining limitations}

\begin{itemize}
\item Cross-sectional, post-mortem, lethal cases only (primary).
\item Replication cohort does not include Delorey (GSE171524) or Wendisch
      (PRJNA705687); these datasets are not on cellxgene census and would require
      raw GEO retrieval.
\item Replication cohort uses a curated 84-gene panel; full-transcriptome HVG/PCA/DPT
      replication is not yet done.
\item Fine program ordering is hypothesis-generating only.
\end{itemize}

\section{What to tell a computer scientist vs a biologist}

\textbf{For the computer scientist:} we built an ablation-first trajectory-analysis
pipeline, pre-registered three success criteria with distinct geometric
interpretations (centroid displacement vs dispersion vs trajectory enrichment),
and found that a partial success (2/3 criteria) with strong out-of-cohort replication
of the two that passed is more informative than a uniform-pass would be. The
methodological takeaways are (i) dispersion beats centroid displacement as a
robustness-collapse metric; (ii) donor-level pseudobulk matters in imbalanced atlases;
(iii) marker-threshold cell-state classifiers don't transfer; (iv) fix your Harmony
wrapper.

\textbf{For the biologist:} in lethal COVID-19, alveolar epithelial cells don't all
move together to a new damaged state; they \emph{fan out} along a graded axis.
This robustness-collapse signature replicates in an independent 618-donor atlas.
The DATP biology is real but needs manifold-based identification, not threshold
classifiers, to cross cohorts cleanly.

\end{document}
